% ============================================
% Assistant Teaching Professor Cover Letter – UMass Lowell
% ============================================

\documentclass[11pt,letterpaper]{letter}

\usepackage{geometry}
\usepackage{microtype}
\usepackage[hidelinks]{hyperref}

\geometry{left=25mm,right=25mm,top=25mm,bottom=25mm}

\signature{Decory Edwards}
\address{
Department of Economics \\
Johns Hopkins University \\
Baltimore, MD 21218 \\
\texttt{dedwar65@jh.edu}
}

\begin{document}

\begin{letter}{Search Committee \\
Department of Economics \\
University of Massachusetts Lowell \\
Lowell, MA \\ USA}

\opening{Dear Members of the Search Committee,}

I am writing to express my interest in the Assistant Teaching Professor position in the Department of Economics at the University of Massachusetts Lowell. I am a Ph.D.\ candidate in Economics at Johns Hopkins University, advised by Professors Christopher D.~Carroll, Nicholas W.~Papageorge, and Stelios Fourakis, and I expect to receive my degree in May 2026. My research fields are macroeconomics, wealth and income dynamics, and computational economics.

My research agenda focuses on understanding the determinants of wealth inequality and the mechanisms through which heterogeneity in returns to wealth shapes aggregate outcomes. In my job-market paper, I estimate the degree of heterogeneity in individual portfolio returns using micro-data from the Survey of Consumer Finances and simulate a structural model in which households face idiosyncratic return risk. I show that allowing for persistent heterogeneity in returns substantially improves the model’s ability to match the empirical distribution of wealth, offering a tractable way to connect micro-level return dispersion to macroeconomic inequality.

In a second project, I use the Health and Retirement Study to examine the relationship between interpersonal trust and portfolio performance. Building on the literature that finds a hump-shaped relationship between trust and income, I show that a similar relationship holds for individual returns to wealth measured in the HRS data, highlighting a behavioral channel through which social attitudes shape saving and investment outcomes. This work also provides new estimates of the persistent component of returns to net worth, which motivate the calibration of heterogeneity in my structural model.

The teaching-focused mission of the University of Massachusetts Lowell strongly aligns with my professional goals. UMass Lowell’s emphasis on student success, applied learning, and inclusive access to higher education makes it an especially attractive environment for faculty who value classroom excellence and sustained student engagement. I am particularly drawn to the opportunity to teach core economics courses that serve diverse student populations and prepare students for careers in business, policy, and public service.

\textbf{Teaching.} I have experience teaching and supporting courses in \textit{Principles of Macroeconomics}, \textit{Intermediate Macroeconomic Theory}, and applied quantitative economics. My teaching emphasizes clarity, economic intuition, and the use of data to connect economic models to real-world outcomes. In a teaching-focused role at UMass Lowell, I would be well prepared to teach introductory and intermediate economics courses and to contribute to curriculum development that supports student learning across majors. I place a strong emphasis on inclusive pedagogy, active learning, and clear communication, particularly for students encountering economics for the first time.

I would be enthusiastic about contributing to the Department of Economics at UMass Lowell through high-quality instruction, student mentoring, and service. Thank you very much for your time and consideration. I would welcome the opportunity to discuss my application further.

\closing{Sincerely,}

\end{letter}
\end{document}
